\documentclass[12pt]{article}
\usepackage[margin=1in]{geometry}
\usepackage{hyperref}
\usepackage{booktabs}
\usepackage{parskip}
\usepackage{enumitem}
\usepackage{titlesec}
\usepackage{xcolor}
\usepackage{mdframed}

\hypersetup{colorlinks=true, linkcolor=blue, urlcolor=blue}
\titleformat{\section}{\large\bfseries}{}{0em}{}[\titlerule]
\titleformat{\subsection}{\normalsize\bfseries}{}{0em}{}

\title{\textbf{OHATS Path Forward}\\
\large Organization Health Assessment Tool \& Survey\\
Archive Overview and Recommendations for Future Research Teams}
\author{Brian Hanley — E483/E583 Information Visualization, Spring 2026\\
Indiana University Luddy School of Informatics, Computing, and Engineering}
\date{April 2026}

\begin{document}
\maketitle
\tableofcontents
\newpage

% -------------------------------------------------------------------------
\section{What Is OHATS}

\textbf{OHATS} (Organization Health Assessment Tool \& Survey) was a web-based
system designed for the \textbf{Tutor/Mentor Connection} (T/MC) by researcher
Steve Roussos between 1998 and 2002. Its purpose was to help youth-serving
nonprofit organizations self-assess their performance against the four
strategic pillars that Daniel F.\ Bassill used to guide the T/MC mission:

\begin{enumerate}[leftmargin=1.5em]
  \item \textbf{Collect information} — document what organizations do and who they reach.
  \item \textbf{Spread information} — share knowledge across the nonprofit network.
  \item \textbf{Build capacity} — develop organizational skills and leadership.
  \item \textbf{Evaluate impact} — measure outcomes and adapt strategy.
\end{enumerate}

OHATS asked organizations to document what they did in relation to one or
more of these pillars. The vision was to aggregate those self-assessments
across the network so that T/MC could see which organizations were active
in which pillars, track change over time, and identify where capacity-building
support was most needed.

The system was active in some form from approximately 2000 to 2002. Rebuild
attempts occurred in 2007 (internal notes) and 2008 (Ganesh Mahalingam,
external developer; SVHATS~2.0 specification by V.~Helde). None of
the rebuilds resulted in a publicly deployed, sustained system.

\begin{mdframed}[backgroundcolor=yellow!10, linecolor=orange!60]
\textbf{Why OHATS matters now (April 2026):} In the April~22,~2026
client meeting, Daniel F.\ Bassill confirmed that OHATS is a priority
for the next semester project. He described the data as directly connected
to his four strategic pillars and said that visualizing OHATS assessment
data would be ``data-visualization-heavy'' in a way that is ideal for the
E483/E583 course. The Spring~2026 team did preliminary cleanup of the
archive; a future team should begin with this document and the folder
inventory below.
\end{mdframed}

% -------------------------------------------------------------------------
\section{Archive Location and Inventory}

All OHATS materials are stored at:
\begin{center}
\texttt{ProjectFiles/OHATS-TMC-SteveRoussos 1998-2002-20260422T193300Z-3-001/\\
OHATS-TMC-SteveRoussos 1998-2002/}
\end{center}

\subsection{Planning Correspondence (1998--2001)}

These documents trace the origin and design of OHATS from Roussos's
initial letter to T/MC through the first operational period.

\begin{center}
\begin{tabular}{lp{9cm}}
\toprule
File & Description \\
\midrule
\texttt{RossosLetter-toTMC-12-6-1998.pdf} & Roussos's initial letter proposing evaluation system to T/MC \\
\texttt{S.Roussos-letter~12-6-1998.pdf} & Companion letter \\
\texttt{S.Roussos-letter-collaboration.ideas-1-22-1999.pdf} & Early collaboration ideas \\
\texttt{S.Roussos-letter-grants.ideas-1999.pdf} & Grant-funding concepts for the system \\
\texttt{S.Roussos-email-.communication-1999.pdf} & 1999 email communications \\
\texttt{S.Roussos-email-.idea-sharing-1999.pdf} & Idea-sharing exchange \\
\texttt{S.Roussos-email-ICQ-net.communication-1999.pdf} & ICQ-based communication \\
\texttt{ConfPresentation-IYP-funding-Roussos-draft-1999.pdf} & Conference presentation, IYP funding \\
\bottomrule
\end{tabular}
\end{center}

\subsection{System Design and Workplans (2000--2002)}

These documents define what OHATS was supposed to do and how it was built.
\textbf{Start here when trying to understand the system schema.}

\begin{center}
\begin{tabular}{lp{9cm}}
\toprule
File & Description \\
\midrule
\texttt{Workplan-to-establish-TMC-OHATS-4-24-2000-Roussos.pdf} & Primary workplan for establishing OHATS \\
\texttt{WorkPlan-TMC-EvalSystem-6-9-2000.pdf} & Evaluation system workplan \\
\texttt{OHATS for TMC-S.Roussos-6-2000.pdf} & Core OHATS description document \\
\texttt{OHATS - what is it.doc} & Plain-language explanation of OHATS \\
\texttt{OHATS development - notes - 2000-01.pdf} & Developer notes during active build \\
\texttt{OHATS-draft-comments-1-7-2001.pdf} & Draft review comments \\
\texttt{OHATS improvement specs - 2-21-2001-Roussos.pdf} & Improvement specifications \\
\texttt{OHATS planning - Roussos email 6-21-2001.pdf} & Planning email \\
\texttt{OHATS - Procedures for Dev \& Application - 3-2002.pdf} & Procedures document \\
\texttt{Mgt-Info-Learning-System-MILS-forCC-7-27-2000-Roussos.pdf} & Broader information system context \\
\texttt{CC-MtInfo-LearningSystem-S.Roussos-7-27-2000.pdf} & Companion MILS document \\
\texttt{MILS for KidsConnection-Draft-8-3-2000-Roussos.pdf} & Application to kids programs \\
\bottomrule
\end{tabular}
\end{center}

\subsection{Reports and Outcomes (2000--2005)}

\begin{center}
\begin{tabular}{lp{9cm}}
\toprule
File & Description \\
\midrule
\texttt{TMC-OHATS Report-March2002.pdf} & Primary operational report \\
\texttt{OHATS-files/TMC\_OHATS\_Report\_Sep-00\_thru\_Mar-02.pdf} & \textbf{Key data source}: full 18-month report Sep 2000--Mar 2002 \\
\texttt{OHATS items reported by G.~Schoen for TMC 2000-2001.pdf} & Items logged by staff \\
\texttt{TMC OHATS development 3-17-2003.pdf} & 2003 development status \\
\texttt{Long-term-impact-CC-TMC-slides-for-discussion~2005.pdf} & 2005 retrospective slides \\
\texttt{TQM chart-circa2000-showing people involved.pdf} & TQM organizational chart \\
\bottomrule
\end{tabular}
\end{center}

\subsection{Funding Documents}

\begin{center}
\begin{tabular}{lp{9cm}}
\toprule
File & Description \\
\midrule
\texttt{Request-for-2ndYr-Funding-OHATS-2001.pdf} & Request for second-year funding \\
\texttt{SRoussos-contract-building-evaluation-OHATS-2000-01.pdf} & Roussos contract \\
\texttt{Research and Information System Associate (RISA)...pdf} & RISA role description for OHATS \\
\texttt{Mentoring Network Grant-abstract-draft~3-1999.pdf} & Grant abstract \\
\texttt{Mentoring Network Grant-problem description-draft~3-1999.pdf} & Grant problem statement \\
\texttt{Workplan for TMCTech-Intern-Summer2001.pdf} & Intern workplan \\
\texttt{Endorsement of S.~Roussos - 2011.pdf} & 2011 endorsement letter \\
\bottomrule
\end{tabular}
\end{center}

\subsection{Rebuild Attempts (2007--2008)}

\begin{center}
\begin{tabular}{lp{9cm}}
\toprule
File & Description \\
\midrule
\texttt{OHATS rebuild - 2-16-2007.pdf} & 2007 rebuild plan \\
\texttt{OHATS rebuild - screen shots - 2007.pdf} & Screen shots of rebuilt version \\
\texttt{OHATS rebuild - screen shots w notes - 2007.pdf} & Annotated screen shots \\
\texttt{OHATS 2007 developer - Ganesh Mahalingam-12-24-2007.pdf} & Developer communication \\
\texttt{SVHATS\_2.0.doc} & \textbf{Most complete requirements document}: SVHATS~2.0 spec \\
\texttt{SVHATS-JourneyMapping-V.Helde-12-4-2008.pdf} & Journey-mapping for 2008 redesign \\
\texttt{OHATS-files/OHATS08report.doc} & 2008 status report \\
\texttt{OHATS-files/OHATS10\_26.doc}, \texttt{OHATS10\_29.doc} & 2008 working documents \\
\texttt{OHATS-files/OHATS~3\_19\_08.doc}, \texttt{OHATS~8\_16\_08.doc} & 2008 meeting notes \\
\texttt{OHATS-files/OHATS\_-\_Advanced\_Features\_-\_Price\_Info[1].pdf} & Pricing for advanced features \\
\texttt{OHATS-files/TMC~Flowchart1.doc} & System flowchart \\
\texttt{OHATS-files/TMC~goals-GISrev.doc} & Goals document with GIS revision \\
\bottomrule
\end{tabular}
\end{center}

\subsection{OHATS Source Code (ASP.NET, circa 2002)}

The \texttt{OHATS/} subfolder contains the original ASP.NET web
application. \textbf{Do not attempt to run this code}---it is legacy
technology (Microsoft FrontPage / .NET 1.x, circa 2002) and serves
only as a documentation artifact showing the system's data schema and
page structure.

Key pages that reveal the data model:
\begin{itemize}[leftmargin=1.5em]
  \item \texttt{Default.aspx} — Home/login
  \item \texttt{EventInfo.aspx}, \texttt{Eventsearchcontent.aspx} — Event tracking
  \item \texttt{Contactinfo.aspx}, \texttt{ContactInfoListcontent.aspx} — Organization contacts
  \item \texttt{Metrics.aspx}, \texttt{AdvancedMetrics.aspx} — Metric reporting
  \item \texttt{CredentialMetrics.aspx}, \texttt{CredentialsAdvancedMetrics.aspx} — Credential tracking
  \item \texttt{Update.aspx}, \texttt{MyContacts.aspx} — Self-service update
\end{itemize}

The page names map directly to the data categories OHATS tracked:
\textbf{events, contacts, metrics, and credentials}---each corresponding
to one or more of Bassill's four strategic pillars.

% -------------------------------------------------------------------------
\section{Data Extraction Opportunities}

\subsection{Where Data Actually Lives}

OHATS was a live web application; the database it wrote to is not included
in this archive. What survives are:

\begin{enumerate}[leftmargin=1.5em]
  \item \textbf{Narrative reports} describing aggregate results
        (\texttt{TMC\_OHATS\_Report\_Sep-00\_thru\_Mar-02.pdf},
        \texttt{TMC-OHATS Report-March2002.pdf}).
  \item \textbf{Working documents} with structured lists of items, people,
        and organizations (\texttt{OHATS items reported by G.~Schoen...pdf},
        \texttt{OHATS development - notes - 2000-01.pdf}).
  \item \textbf{The source code schema} (visible from \texttt{.aspx} page names
        and the \texttt{SVHATS\_2.0.doc} requirements).
  \item \textbf{The 2007--2008 rebuild screen shots}, which show actual data
        populated in the system's UI during those rebuild periods.
\end{enumerate}

\subsection{Recommended Extraction Steps for a Future Team}

\begin{enumerate}[leftmargin=1.5em]
  \item \textbf{Read the 18-month report first.}
        \texttt{OHATS-files/TMC\_OHATS\_Report\_Sep-00\_thru\_Mar-02.pdf}
        is the primary record of what the system actually captured. Extract
        any tabular data (organization names, pillars addressed, assessment
        scores) into a CSV for analysis.

  \item \textbf{Read SVHATS\_2.0.doc.}
        This is the most complete requirements specification. It defines
        the data model that the 2008 rebuild was targeting. Use it to design
        a modern schema that could receive OHATS-style data.

  \item \textbf{Review the 2007 screen shots with notes.}
        \texttt{OHATS rebuild - screen shots w notes - 2007.pdf} shows the
        UI with annotated data. This may contain organization names and
        assessment values that can be extracted manually.

  \item \textbf{Cross-reference with conference attendance data.}
        The conference attendance dataset (6,410 records, 2,163 unique
        organizations) covers the same 1994--2014 period. Joining
        OHATS-assessed organizations to the conference attendance data
        would reveal which organizations both attended conferences
        \emph{and} participated in self-assessment---a proxy for
        ``most engaged'' network members.

  \item \textbf{Design the data entry tool.}
        The long-term goal is a modern replacement for OHATS: a Django
        (or similar) form that lets organizations self-report against
        Bassill's four pillars. The \texttt{MappingAndAnalyzingParticipation}
        and \texttt{YourConferenceNetwork} codebases provide a technical
        starting point.
\end{enumerate}

% -------------------------------------------------------------------------
\section{Recommended Visualization Directions}

Once structured OHATS data is extracted or collected, the following
visualizations would be high-impact for Bassill's stakeholders:

\begin{enumerate}[leftmargin=1.5em]
  \item \textbf{Pillar coverage heatmap.}
        Rows = organizations; columns = the four strategic pillars.
        Color encodes self-reported engagement level (e.g., never, partially,
        consistently). This directly answers ``who is doing what?''

  \item \textbf{Pillar distribution by SNA sector.}
        Do Program-sector organizations prioritize different pillars than
        Foundation or Business orgs? A stacked bar or grouped chart
        answers this and connects OHATS findings to the existing conference
        attendance SNA categories.

  \item \textbf{Timeline of OHATS use.}
        A simple time-series showing which organizations reported in which
        years (2000--2002, and if data exists from 2007--2008) would
        show engagement patterns over the system's lifetime.

  \item \textbf{OHATS $\times$ conference co-attendance network.}
        Overlay OHATS participation onto the org-org co-attendance network
        already built in Kumu.io. Organizations that both attended conferences
        \emph{and} self-assessed become a highlighted subset of the full
        network---visualizing the most engaged segment.

  \item \textbf{Radar/spider chart per organization.}
        For individual organizations that completed OHATS, a radar chart
        showing scores across the four pillars gives Bassill a one-page
        org profile.
\end{enumerate}

% -------------------------------------------------------------------------
\section{Technical Handoff Notes}

\begin{itemize}[leftmargin=1.5em]
  \item The ASP.NET source code in \texttt{OHATS/} is informational only.
        Do not attempt to run, compile, or modernize it.
  \item The \texttt{SVHATS\_2.0.doc} and \texttt{OHATS-files/TMC~Flowchart1.doc}
        are the best references for system design intent.
  \item Any new OHATS data entry should be confirmed with Bassill
        before development. He has strong views on category definitions
        and pillar assignments (see April~22,~2026 transcript in
        \texttt{ProjectFiles/transcript\_video2125685008.txt}).
  \item The existing Django codebase (\texttt{MappingAndAnalyzingParticipation/})
        can be extended with a new Django app for OHATS data---the ORM,
        \texttt{load\_data} pattern, and export endpoints are all reusable.
  \item Connect with Bassill via \texttt{tutormentor1@gmail.com} and review
        the ``Strategy'' concept map on
        \href{https://tutormentorexchange.net}{tutormentorexchange.net}
        (under Tutor/Mentor Connection $\to$ Strategy) before beginning
        any OHATS visualization work.
  \item Review the full entry in \texttt{ProjectFiles/future\_team\_suggestions.tex}
        (Section~4.1, ``OHATS data integration'') for additional context.
\end{itemize}

% -------------------------------------------------------------------------
\section{Connection to the Broader Project}

OHATS is not a standalone project---it is the next layer of the same
ecosystem the Spring~2026 team built. The four-pillar strategy that OHATS
tracked is visible in the Kumu concept maps Bassill showed in the April~22
meeting. The conference attendance data already documents \emph{who} was
involved; OHATS would document \emph{what they did} and \emph{how well}.
Together, they answer the question Bassill has been asking for three decades:
who needs to be in the network, and are they doing the work that makes the
network effective?

A future team that connects these two datasets---attendance + assessment---will
have produced something qualitatively more powerful than either alone.

\end{document}
